\begin{lemma}
\label{lemma-more-base-change}
In diagram (\ref{equation-base-change}) assume
\begin{enumerate}
\item $g : Y' \to Y$ is a morphism of affine schemes,
\item $f : X \to Y$ is proper, and
\item $f$ and $g$ are Tor independent.
\end{enumerate}
Then the base change map (\ref{equation-base-change-map}) induces an
isomorphism
$$
L(g')^*a(K) \longrightarrow a'(Lg^*K)
$$
in the following cases
\begin{enumerate}
\item for all $K \in D_\QCoh(\mathcal{O}_Y)$ if $f$
is flat of finite presentation,
\item for all $K \in D_\QCoh(\mathcal{O}_Y)$ if $f$
is perfect and $Y$ Noetherian,
\item for $K \in D_\QCoh^+(\mathcal{O}_Y)$ if $g$ has finite Tor dimension
and $Y$ Noetherian.
\end{enumerate}
\end{lemma}

\begin{proof}
Write $Y = \Spec(A)$ and $Y' = \Spec(A')$. As a base change of an affine
morphism, the morphism $g'$ is affine. Let $M$ be a perfect generator
for $D_\QCoh(\mathcal{O}_X)$, see Derived Categories of Schemes, Theorem
\ref{perfect-theorem-bondal-van-den-Bergh}. Then $L(g')^*M$ is a
generator for $D_\QCoh(\mathcal{O}_{X'})$, see
Derived Categories of Schemes, Remark \ref{perfect-remark-pullback-generator}.
Hence it suffices to show that (\ref{equation-base-change-map})
induces an isomorphism
\begin{equation}
\label{equation-iso}
R\Hom_{X'}(L(g')^*M, L(g')^*a(K))
\longrightarrow
R\Hom_{X'}(L(g')^*M, a'(Lg^*K))
\end{equation}
of global hom complexes, see
Cohomology, Section \ref{cohomology-section-global-RHom},
as this will imply the cone of $L(g')^*a(K) \to a'(Lg^*K)$
is zero.
The structure of the proof is as follows: we will first show that
these Hom complexes are isomorphic and in the last part of the proof
we will show that the isomorphism is induced by (\ref{equation-iso}).

\medskip\noindent
The left hand side. Because $M$ is perfect, the canonical map
$$
R\Hom_X(M, a(K)) \otimes^\mathbf{L}_A A'
\longrightarrow
R\Hom_{X'}(L(g')^*M, L(g')^*a(K))
$$
is an isomorphism by Derived Categories of Schemes, Lemma
\ref{perfect-lemma-affine-morphism-and-hom-out-of-perfect}.
We can combine this with the isomorphism
$R\Hom_Y(Rf_*M, K) = R\Hom_X(M, a(K))$
of Lemma \ref{lemma-iso-global-hom}
to get that the left hand side equals
$R\Hom_Y(Rf_*M, K) \otimes^\mathbf{L}_A A'$.

\medskip\noindent
The right hand side. Here we first use the isomorphism
$$
R\Hom_{X'}(L(g')^*M, a'(Lg^*K)) = R\Hom_{Y'}(Rf'_*L(g')^*M, Lg^*K)
$$
of Lemma \ref{lemma-iso-global-hom}. Then we use the base change
map $Lg^*Rf_*M \to Rf'_*L(g')^*M$ is an isomorphism by
Derived Categories of Schemes, Lemma \ref{perfect-lemma-compare-base-change}.
Hence we may rewrite this as $R\Hom_{Y'}(Lg^*Rf_*M, Lg^*K)$.
Since $Y$, $Y'$ are affine and $K$, $Rf_*M$ are in $D_\QCoh(\mathcal{O}_Y)$
(Derived Categories of Schemes, Lemma
\ref{perfect-lemma-quasi-coherence-direct-image})
we have a canonical map
$$
\beta :
R\Hom_Y(Rf_*M, K) \otimes^\mathbf{L}_A A'
\longrightarrow
R\Hom_{Y'}(Lg^*Rf_*M, Lg^*K)
$$
in $D(A')$. This is the arrow
More on Algebra, Equation (\ref{more-algebra-equation-base-change-RHom})
where we have used Derived Categories of Schemes, Lemmas
\ref{perfect-lemma-affine-compare-bounded} and
\ref{perfect-lemma-quasi-coherence-internal-hom}
to translate back and forth into algebra.
\begin{enumerate}
\item If $f$ is flat and of finite presentation, the complex $Rf_*M$
is perfect on $Y$ by Derived Categories of Schemes, Lemma
\ref{perfect-lemma-flat-proper-perfect-direct-image-general}
and $\beta$ is an isomorphism by
More on Algebra, Lemma \ref{more-algebra-lemma-base-change-RHom} part (1).
\item If $f$ is perfect and $Y$ Noetherian, the complex $Rf_*M$
is perfect on $Y$ by More on Morphisms, Lemma
\ref{more-morphisms-lemma-perfect-proper-perfect-direct-image}
and $\beta$ is an isomorphism as before.
\item If $g$ has finite tor dimension and $Y$ is Noetherian,
the complex $Rf_*M$ is pseudo-coherent on $Y$
(Derived Categories of Schemes, Lemmas
\ref{perfect-lemma-direct-image-coherent} and
\ref{perfect-lemma-identify-pseudo-coherent-noetherian})
and $\beta$ is an isomorphism by
More on Algebra, Lemma \ref{more-algebra-lemma-base-change-RHom} part (4).
\end{enumerate}
We conclude that we obtain the same answer as in the previous paragraph.

\medskip\noindent
In the rest of the proof we show that the identifications of
the left and right hand side of (\ref{equation-iso})
given in the second and third paragraph are in fact given by
(\ref{equation-iso}). To make our formulas manageable
we will use $(-, -)_X = R\Hom_X(-, -)$, use $- \otimes A'$
in stead of $- \otimes_A^\mathbf{L} A'$, and we will abbreviate
$g^* = Lg^*$ and $f_* = Rf_*$. Consider the following
commutative diagram
$$
\xymatrix{
((g')^*M, (g')^*a(K))_{X'} \ar[d] &
(M, a(K))_X \otimes A' \ar[l]^-\alpha \ar[d] &
(f_*M, K)_Y \otimes A' \ar@{=}[l] \ar[d] \\
((g')^*M, (g')^*a(g_*g^*K))_{X'} &
(M, a(g_*g^*K))_X \otimes A' \ar[l]^-\alpha &
(f_*M, g_*g^*K)_Y \otimes A' \ar@{=}[l] \ar@/_4pc/[dd]_{\mu'} \\
((g')^*M, (g')^*g'_*a'(g^*K))_{X'} \ar[u] \ar[d] &
(M, g'_*a'(g^*K))_X \otimes A' \ar[u] \ar[l]^-\alpha \ar[ld]^\mu &
(f_*M, K) \otimes A' \ar[d]^\beta \\
((g')^*M, a'(g^*K))_{X'} &
(f'_*(g')^*M, g^*K)_{Y'} \ar@{=}[l] \ar[r] &
(g^*f_*M, g^*K)_{Y'}
}
$$
The arrows labeled $\alpha$ are the maps from
Derived Categories of Schemes, Lemma
\ref{perfect-lemma-affine-morphism-and-hom-out-of-perfect}
for the diagram with corners $X', X, Y', Y$.
The upper part of the diagram is commutative as the horizontal arrows are
functorial in the entries.
The middle vertical arrows come from the invertible transformation
$g'_* \circ a' \to a \circ g_*$  of Lemma \ref{lemma-flat-precompose-pus}
and therefore the middle square is commutative.
Going down the left hand side is (\ref{equation-iso}).
The upper horizontal arrows provide the identifications used in the
second paragraph of the proof.
The lower horizontal arrows including $\beta$ provide the identifications
used in the third paragraph of the proof. Given $E \in D(A)$,
$E' \in D(A')$, and $c : E \to E'$ in $D(A)$ we will denote
$\mu_c : E \otimes A' \to E'$ the map induced by $c$
and the adjointness of restriction and base change;
if $c$ is clear we write $\mu = \mu_c$, i.e., we
drop $c$ from the notation. The map $\mu$ in the diagram is of this
form with $c$ given by the identification
$(M, g'_*a(g^*K))_X = ((g')^*M, a'(g^*K))_{X'}$
; the triangle involving $\mu$ is commutative by
Derived Categories of Schemes, Remark \ref{perfect-remark-multiplication-map}.

\medskip\noindent
Observe that
$$
\xymatrix{
(M, a(g_*g^*K))_X &
(f_*M, g_* g^*K)_Y \ar@{=}[l] &
(g^*f_*M, g^*K)_{Y'} \ar@{=}[l] \\
(M, g'_* a'(g^*K))_X \ar[u] &
((g')^*M, a'(g^*K))_{X'} \ar@{=}[l] &
(f'_*(g')^*M, g^*K)_{Y'} \ar@{=}[l] \ar[u]
}
$$
is commutative by the very definition of the transformation
$g'_* \circ a' \to a \circ g_*$. Letting $\mu'$ be as above
corresponding to the identification
$(f_*M, g_*g^*K)_X = (g^*f_*M, g^*K)_{Y'}$, then the
hexagon commutes as well. Thus it suffices to show that
$\beta$ is equal to the composition of
$(f_*M, K)_Y \otimes A' \to (f_*M, g_*g^*K)_X \otimes A'$
and $\mu'$. To do this, it suffices to prove the two induced maps
$(f_*M, K)_Y \to (g^*f_*M, g^*K)_{Y'}$ are the same.
In other words, it suffices to show the diagram
$$
\xymatrix{
R\Hom_A(E, K) \ar[rr]_{\text{induced by }\beta} \ar[rd] & &
R\Hom_{A'}(E \otimes_A^\mathbf{L} A', K \otimes_A^\mathbf{L} A') \\
& R\Hom_A(E, K \otimes_A^\mathbf{L} A') \ar[ru]
}
$$
commutes for all $E, K \in D(A)$. Since this is how $\beta$ is constructed in
More on Algebra, Section \ref{more-algebra-section-base-change-RHom}
the proof is complete.
\end{proof}